% ------------------------------------------------------------
% Python Diagnostic Toolkit (Module 5)
% Author: Borealis S. Hedling
% Altitude: A6.5 (Boundary Layer)
% Purpose: Provide Python-style pseudocode for runtime
%          diagnostics corresponding to the A6 firewall
%          invariants. These tools allow curvature, spectral
%          mismatch, holonomy drift, collapse locality, and
%          stability envelopes to be evaluated without requiring
%          non-dual algebra or projection-repair operators.
% ------------------------------------------------------------

\section{Python Diagnostic Toolkit}

This module provides Python-style pseudocode for evaluating the A6 firewall
invariants at runtime. These diagnostic tools mirror the Lean verification
lemmas and allow the Boundary Mapper to compute curvature deviation, spectral
mismatch, holonomy drift, collapse locality, and stability envelope
monotonicity.

\subsection{Toolkit Philosophy}

The diagnostic toolkit is intentionally algebra-light. It evaluates geometric
and spectral invariants without requiring access to A7 non-dual algebra or
projection-repair operators. This ensures that diagnostics remain readable
and accessible.

\section{Curvature Estimator}

\begin{verbatim}
def curvature_deviation(K, p, q):
    """
    Estimate curvature deviation between adjacent points.
    K : function returning curvature at a point
    p, q : adjacent points
    """
    return abs(K(p) - K(q))
\end{verbatim}

This function evaluates curvature continuity across the firewall boundary.

\section{Spectral Mismatch Calculator}

\begin{verbatim}
def spectral_mismatch(L, p, q):
    """
    Compute spectral mismatch between operator spectra.
    L : function returning operator at a point
    p, q : adjacent points
    """
    spec_p = spectrum(L(p))
    spec_q = spectrum(L(q))
    return distance(spec_p, spec_q)
\end{verbatim}

This function detects spectral tearing or instability.

\section{Holonomy Drift Calculator}

\begin{verbatim}
def holonomy_drift(H, gamma):
    """
    Compute holonomy drift along a loop gamma.
    H : function returning holonomy of a loop
    gamma : loop in the manifold
    """
    return H(gamma)
\end{verbatim}

This function evaluates drift accumulation and checks fracture thresholds.

\section{Collapse-Surface Detector}

\begin{verbatim}
def collapse_locality(Sigma, p, eps):
    """
    Check whether collapse surfaces are locally detectable.
    Sigma : set of collapse-surface points
    p : point in the manifold
    eps : locality radius
    """
    neighborhood = ball(p, eps)
    return len(neighborhood.intersection(Sigma)) > 0
\end{verbatim}

This function ensures collapse surfaces do not smear across the boundary.

\section{Stability Envelope Evaluator}

\begin{verbatim}
def stability_monotonicity(E, p, q):
    """
    Evaluate stability envelope monotonicity.
    E : function returning envelope value
    p, q : points with E(p) <= E(q)
    """
    if E(p) <= E(q):
        return True
    return False
\end{verbatim}

This function checks whether stability envelopes remain monotone.

\section{Toolkit Summary}

The Python diagnostic toolkit provides runtime evaluation of the A6 firewall
invariants. These tools allow the Boundary Mapper to compute curvature
deviation, spectral mismatch, holonomy drift, collapse locality, and
stability envelope monotonicity without requiring non-dual algebra. The
toolkit ensures that diagnostics remain accessible, modular, and safe.
